\documentclass[12pt,a4paper]{article}

\usepackage{fontspec}
\setmainfont{Verdana}
\usepackage[french]{babel}
\usepackage{geometry}
\usepackage{setspace}
\usepackage{array}
\usepackage{enumitem}
\usepackage{amsmath}
\usepackage{tikz}

\geometry{margin=2cm}
\onehalfspacing
\pagestyle{empty}

\setlength{\parindent}{0pt}
\setlist[enumerate]{label=\textbf{\alph*.}, leftmargin=1.2cm}

\newcommand{\tableauproportion}[4]{
\begin{tikzpicture}[baseline={(current bounding box.center)}]
    \draw (0,0) rectangle (2.4,1.4);
    \draw (1.2,0) -- (1.2,1.4);
    \draw (0,0.7) -- (2.4,0.7);

    \node at (0.6,1.05) {#1};
    \node at (1.8,1.05) {#2};
    \node at (0.6,0.35) {#3};
    \node at (1.8,0.35) {#4};

    \draw[->, thick] (-0.35,1.1) to[out=200,in=160] (-0.35,0.3);
    \node[draw, inner sep=3pt] at (-1.15,0.7) {$\times \ldots$};

    \draw[->, thick] (2.75,0.3) to[out=-20,in=20] (2.75,1.1);
    \node[draw, inner sep=3pt] at (3.55,0.7) {$\div \ldots$};
\end{tikzpicture}
}

\newcommand{\Entete}[1]{
\begin{center}
    {\Large \textbf{Contrôle séquence 13}}\\
    \textbf{Version #1}\\
\vspace{1cm}
    Calculatrice autorisée.

    Chaque réponse doit être justifiée par un calcul.
\end{center}
}

\newcommand{\ExercicePommes}[5]{
\vspace{0.5cm}

\underline{Exercice 1}

Avec #1 kg de pommes, on peut faire #2 L de jus de pomme.

\textbf{Recopier} et compléter le tableau de proportionnalité suivant.

\vspace{0.3cm}

\begin{center}
\renewcommand{\arraystretch}{1.5}
\begin{tabular}{|>{\raggedright\arraybackslash}m{5cm}|c|c|c|c|c|}
\hline
Masse de pommes en kg & #1 & & #4 & #5 & 1 \\
\hline
Volume de jus de pomme en L & & #3 & & & \\
\hline
\end{tabular}
\end{center}
}

\newcommand{\ExerciceAmanda}[2]{
\vspace{0.5cm}

\underline{Exercice 2}

Le véhicule d'Amanda consomme en moyenne #1 L sur un trajet de 100 km.

Calculer la distance que peut espérer parcourir Amanda avec #2 L d'essence.
}

\newcommand{\ExerciceCoef}[8]{
\vspace{0.5cm}

\underline{Exercice 3}

Les tableaux ci-dessous sont des tableaux de proportionnalité. \textbf{Recopie} les tableaux et donne le coefficient de proportionnalité à écrire à la place des pointillés.

\vspace{0.5cm}

\begin{center}
\tableauproportion{#1}{#2}{#3}{#4}
\qquad\qquad
\tableauproportion{#5}{#6}{#7}{#8}
\end{center}
}

\newcommand{\ExercicePieton}[4]{
\vspace{0.5cm}

\underline{Exercice 4}

Un piéton marche à une allure régulière. Il parcourt #1 km en #2 min.

\begin{enumerate}
    \item Quelle distance, exprimée en km, parcourt-il en #3 minutes ?
    \item Quel temps met-il à parcourir #4 km ?
\end{enumerate}
}

\newcommand{\ExerciceYaourts}[4]{
\vspace{6 cm}

\underline{Exercice 5}

On considère les prix suivants :

\vspace{0.3cm}

\begin{center}
\renewcommand{\arraystretch}{1.5}
\begin{tabular}{|>{\centering\arraybackslash}m{5cm}|>{\centering\arraybackslash}m{5cm}|}
\hline
\textbf{Nombre de yaourts} & \textbf{Prix} \\
\hline
#1 yaourts & $#2$ € \\
\hline
#3 yaourts & $#4$ € \\
\hline
\end{tabular}
\end{center}

\vspace{0.3cm}

Le prix de ces yaourts est-il proportionnel au nombre de yaourts ? Expliquer pourquoi.
}

\newcommand{\ExerciceAthlete}[4]{
\vspace{0.5cm}

\underline{Exercice 6}

Un athlète court le #1 m en #2 s et le #3 m en #4 secondes.

La distance parcourue par cet athlète est-elle proportionnelle au temps de parcours ? Justifie ta réponse.
}

\newcommand{\VersionControle}[1]{
\Entete{#1}
}

\begin{document}

\VersionControle{1}
\ExercicePommes{100}{60}{6}{20}{130}
\ExerciceAmanda{5}{32}
\ExerciceCoef{3}{5}{12}{20}{7}{56}{9}{72}
\ExercicePieton{$5{,}1$}{60}{48}{7}
\ExerciceYaourts{4}{1{,}20}{8}{2{,}40}
\ExerciceAthlete{100}{10}{200}{22}
\clearpage

\VersionControle{2}
\ExercicePommes{80}{48}{12}{30}{130}
\ExerciceAmanda{10}{42}
\ExerciceCoef{4}{6}{20}{30}{8}{64}{10}{80}
\ExercicePieton{$4{,}8$}{50}{40}{$6{,}4$}
\ExerciceYaourts{3}{0{,}90}{6}{1{,}80}
\ExerciceAthlete{80}{8}{160}{18}
\clearpage

\VersionControle{3}
\ExercicePommes{120}{72}{18}{40}{190}
\ExerciceAmanda{2}{49}
\ExerciceCoef{2}{7}{6}{21}{9}{45}{12}{60}
\ExercicePieton{$6{,}3$}{70}{56}{$8{,}1$}
\ExerciceYaourts{5}{1{,}75}{10}{3{,}50}
\ExerciceAthlete{50}{5}{100}{11}
\clearpage

\VersionControle{4}
\ExercicePommes{90}{54}{24}{50}{180}
\ExerciceAmanda{20}{56}
\ExerciceCoef{5}{8}{20}{32}{6}{54}{8}{72}
\ExercicePieton{$5{,}6$}{80}{64}{$7{,}7$}
\ExerciceYaourts{6}{2{,}10}{12}{4{,}20}
\ExerciceAthlete{120}{12}{240}{27}
\clearpage

\VersionControle{5}
\ExercicePommes{150}{90}{30}{60}{260}
\ExerciceAmanda{5}{63}
\ExerciceCoef{6}{9}{30}{45}{10}{50}{12}{60}
\ExercicePieton{$4{,}5$}{45}{36}{$6{,}8$}
\ExerciceYaourts{4}{1{,}40}{8}{2{,}80}
\ExerciceAthlete{60}{6}{120}{13}
\clearpage

\VersionControle{6}
\ExercicePommes{70}{42}{36}{80}{210}
\ExerciceAmanda{10}{28}
\ExerciceCoef{4}{7}{16}{28}{8}{40}{11}{55}
\ExercicePieton{$6{,}6$}{55}{44}{$8{,}4$}
\ExerciceYaourts{5}{2{,}25}{10}{4{,}50}
\ExerciceAthlete{90}{9}{180}{20}
\clearpage

\VersionControle{7}
\ExercicePommes{110}{66}{42}{90}{270}
\ExerciceAmanda{2}{65}
\ExerciceCoef{3}{8}{15}{40}{12}{48}{15}{60}
\ExercicePieton{$5{,}4$}{75}{60}{$7{,}2$}
\ExerciceYaourts{3}{1{,}05}{6}{2{,}10}
\ExerciceAthlete{70}{7}{140}{16}
\clearpage

\VersionControle{8}
\ExercicePommes{130}{78}{48}{100}{310}
\ExerciceAmanda{50}{77}
\ExerciceCoef{7}{9}{21}{27}{14}{42}{18}{54}
\ExercicePieton{$7{,}2$}{90}{72}{$9{,}6$}
\ExerciceYaourts{6}{2{,}40}{12}{4{,}80}
\ExerciceAthlete{110}{11}{220}{25}

\end{document}